% Main free-coordinate construction; source-specific metrics may be redefined by assembler.
\providecommand{\FreeCMainDegree}{1\,629\,856}
\providecommand{\FreeCMainOps}{517}
\providecommand{\FreeCParameterDegree}{16\,976}
\section{A free Lucas coordinate: removing the high-degree index carrier}
\label{sec:free-lucas-coordinate}

The central new construction has ten unrestricted integer unknowns.
Its literal source-specific circuit has exact degree $\FreeCMainDegree$ and
$\FreeCMainOps$ live binary arithmetic operations. The mathematical change
is not a better addition chain or a different source layout. We prove that
a high-degree defining congruence in the Lucas certificate is unnecessary
on a fixed residue slice. Its former quotient coordinate becomes a direct,
low-degree Lucas coordinate. The degree formula changes from
$196a+260$ to $96a+160$, where $a=\deg A$.

All statements here concern ordinary mathematics formalizable in ZFC.
The native compiler/refutation equivalence, Pell classification and the
identified Lucas index lemmas remain explicit dependencies. This is not a
new Lean or Rocq certificate, a consistency decision, or a global minimality
claim.

\subsection{The literal source prefix}
The selected positive-gap shared-tap certificate has $W=116$, $H=1873$,
31 weights $w_i$, 22 targets $q_j$, and three kernel polynomials $D_j$.
The complete finite lists are in \texttt{norm/final10/data/}; their
coefficient-extraction identities are proved in the source section below.
Put $\kappa=2134940$, $b=6t+4$, $g=bg_0$, and define
\begin{gather*}
 B=2^\kappa b^2,\quad \ell=\sum_{i=1}^{31}B^{w_i},\quad
 U_c=\ell+g,\quad G_c=D_2(B)U_c^2+D_1(B)U_c+D_0(B),\\
 A_T=\sum_{j=1}^{22}B^{q_j},\quad N_0=B^{117},\quad N=B^{1991},\\
 \mathrm{mask}=N_0-1-(b-1)\ell,\quad J_c=2G_c+A_T,\\
 S_c=g+N_0J_c,\quad T_c=\mathrm{mask}+N_0(B-2)A_T-2,\\
 X=N(S_c+T_c+1)-S_c-1,\\
 Y=2^{16\kappa-23}bB^{2114},\qquad\mathcal C=bB^{116}.
\end{gather*}
Every displayed sum has a fixed finite number of terms, and every exponent
is a fixed integer. The whole native four-port tuple and the twelve
materialized circuit leaves are bound independently; they must not be
substituted by the different source34 tuple used in $P_{28}$ and $P_{26}$.
The new degree bounds are $\deg X=7962$, $\deg Y=4229$ and
$\deg A=16424$. The independently checked primary circuit uses 301
multiplications, 155 additions and 25 subtractions, totaling 481.
Compared with the incoming ten-integer degree-$3{,}327{,}556$/519-operation
construction, both metrics decrease. It also improves the previous
iteration's 494-operation, degree-$3{,}507{,}876$ variant.
The new literal circuit treats the materialized threshold integer
$2^{34159017}$ as one fixed leaf. Retaining the predecessor's five
constant-building gates would give 486 operations, not 481. Both counts
are reported to separate structural simplification from constant folding;
neither is a bit-complexity or execution-time estimate.

\subsection{The complete finalizer}

Take the fixed source polynomials $X,Y,\mathcal C$ and its normalizations
\[
 f_{\rm old}=2t+1,\qquad b=6t+4,\qquad
 g=bg_0,\qquad k=b(3k_0+2).
\]
The exact source prefix and all four coefficient ports are specified in the
accompanying source section and literal data. Its universally valid
congruences are
\[
 X\equiv1\pmod2,\quad X\equiv0\pmod3,\quad
 b\mid X+1,\quad b\mid Y.
\]
The ten supplied variables are
$(t,g_0,c_0,k_0,l,w,x,y,z,m)\in\mathbb Z^{10}$. Every other symbol below
is a polynomial abbreviation. Define
\begin{gather*}
 L=lY,\quad U=2LX,\quad V=4gwY,\quad W=bw,\\
 A_{1/2}=(LX)(V+1),\quad A=2A_{1/2},\quad B_1=2X+1,
 \qquad \boxed{C=6c_0-1},\\
 P_{1/2}=(LX)UV,\quad P=2P_{1/2},\quad
 K=X+1+k(P-2),\\
 S=2A-5,\quad T=3WC-2(W^2-1),\quad
 Q=(P_{1/2}^2-1)K^2+1,\\
 \Delta=8\mathcal C^3gK^2-g^2
       \bigl(32(C-KL)^2\mathcal C^3+g^2K^2\bigr),\\
 \xi=zx,\quad D=(A^2-4)C^2+4,\quad E=C^2D\xi,\\
 F=4(A^2-4)E^2+1,\quad
 G=2(A+2)(A-2)^2E^2-1,\\
 \Omega=16E,\quad H=C+F(B_1-C+\Omega y),\quad
 I=(G^2-1)H^2+1.
\end{gather*}
The circuit uses the identical expression
$G=(A_{1/2}-1)(F-1)-1$ to share operations. The subscripts $1/2$ name
explicit polynomial registers, not divisions in the equation.

Keep the original one-modulus intersective finalizer:
\begin{gather*}
 \rho=T+Sy,\qquad \Sigma=\xi S^5,\\
 \Theta=(\rho^3-\rho S^2-S^3)(\rho^2+23S^2),\\
 r=bQ,\qquad \delta=9\Delta-3,\qquad
 a_0=r+\delta(rz)^2,\quad a_1=DFI,\quad Z=\Sigma m-\Theta.
\end{gather*}
The single equation is
\begin{equation}\label{eq:free-C-final}
 \boxed{\bigl(Z^2-\Sigma^2(a_0+a_1)\bigr)^2
             -4\Sigma^4a_0a_1=0.}
\end{equation}
Every exponent is fixed. There are no hidden quotient, root, binomial,
variable-exponent, or bounded-universal operations. All fixed leaves are
materialized integers with hashes. In particular the new parameterization
has no requirement $(A-2)\mid C-B_1$ built into it.

\subsection{A strengthened signed Lucas certificate}

\paragraph{Lemma.}
Suppose $A,B,C,\xi,y\in\mathbb Z$ satisfy
\[
 6\mid A,\quad B\equiv3\pmod4,\quad
 1<B<|A|/2-1,\quad C\equiv5\pmod6,\quad\xi\ne0.
\]
Use $D,E,F,G,\Omega,H,I$ as above with $B$ in place of $B_1$.
If $DFI$ is a square, then $C=u_B(A)$, where
$u_0=0,u_1=1,u_{n+2}=Au_{n+1}-u_n$.
No hypothesis $(A-2)\mid C-B$ is needed.

\paragraph{Proof.}
Here $C$ is odd and nonzero. Since $|A|>2$, the three integers $D,F,I$
are positive. The certificate identities
\[
 F\equiv I\equiv1\pmod D,\qquad 4I\equiv D\pmod F
\]
and oddness of $F$ make them pairwise coprime, so they are separately
squares. The usual Pell classification gives integers $s,\mu,t$ with
\[
 C=u_s(A),\qquad 4E=u_\mu(A),\qquad H=u_t(2G).
\]
Put $n=|\mu|>3$. The parity identities show that $n$ is even and $s,t$
are odd. The Lucas index-separation lemma applied to the congruence modulo
$F$ gives $s\equiv\pm t\pmod{2n}$. The rank-divisibility lemma applied
to $C^2\mid u_n(A)$ gives $C\mid n$. Finally $H\equiv B\pmod\Omega$
and $2G\equiv-2\pmod\Omega$ give $t\equiv B\pmod\Omega$.
Thus
\begin{equation}\label{eq:free-C-rank-residue}
 s\equiv B\pmod{2C}\quad\hbox{or}\quad s\equiv-B\pmod{2C}.
\end{equation}
These are precisely the congruence-free classification and index steps of
the modified certificate; the new work is to remove its exceptional ranks
and sign without the old defining congruence.

For every $j\ge1$, direct paired induction in the recurrence proves
\begin{equation}\label{eq:free-C-paired-recurrence}
 u_{2j}(A)\equiv(-1)^{j-1}jA\pmod{A^2},\qquad
 u_{2j+1}(A)\equiv(-1)^j\pmod{A^2}.
\end{equation}
The base is $u_2=A,u_3=A^2-1$. The next two recurrence steps multiply the
second congruence by $A$ and subtract the first, then multiply again by $A$
and subtract the second. This works for negative $A$ as well.

If $s=\pm1$, then $C=\pm1$ and $D=A^2$. Therefore $A^2\mid u_n(A)$.
Writing $n=2j$, the first congruence in
\eqref{eq:free-C-paired-recurrence} yields $A\mid j$.
Consequently $2|A|$ divides both $2n$ and
$\Omega=16A^2\xi$. Combining $s\equiv\pm t\pmod{2n}$ with
$t\equiv B\pmod\Omega$ forces $B\equiv\pm1\pmod{2|A|}$,
contradicting $1<B<|A|/2-1$.
This argument does not assume the generally false divisibility $n\mid u_n(A)$.

Hence $|s|\ge2$. Lucas growth gives
$|C|=u_{|s|}(|A|)\ge|A|>B$ and $|s|\le|C|$.
The two congruences in \eqref{eq:free-C-rank-residue} therefore force
$s=B$ or $s=-B$: a nonzero multiple of $2C$ would be too large.
Since $B\equiv3\pmod4$, the second congruence in
\eqref{eq:free-C-paired-recurrence} gives $u_B(A)\equiv-1\pmod{A^2}$,
and thus modulo3. The branch $s=-B$ would give $C\equiv1\pmod3$,
contrary to $C\equiv5\pmod6$. Therefore $s=B$, as required.
\hfill$\square$

The odd, nonzero residue slice is essential: $A=12$, $B=3$, $C=0$, $\xi=1$, $y=0$ gives $DFI=4$, although $u_3(12)=143$.

\paragraph{A reusable unsigned interface.}
The mechanism is not restricted to this particular circuit. For even $A$,
odd $C$, odd $B$ with $1<B<|A|/2-1$, and nonzero $\xi$, the same
square-certificate index argument gives $C=\pm u_B(A)$ without assuming
$A-2\mid C-B$. To handle both exceptional indices $s=\pm1$, use
$C^2=1$: then $D=A^2$, $4E=4A^2\xi$ and the paired recurrence forces
$2|A|\mid n$, just as above. Both exceptions contradict the range of $B$.
For $|s|\ge3$, the growth and index congruences give $s=\pm B$.
The fixed residue conditions $3\mid A$, $B\equiv3\pmod4$,
$C\equiv-1\pmod3$ then select the positive index. Thus the primary
construction is a corollary of a stronger arithmetic interface, not merely
a shorter evaluation DAG. Its dependency on the standard Lucas index
comparison is unchanged; no claim of literature priority is made.

\subsection{Signed soundness of the ordinary equation}

The source and $k$ normalization imply $b\mid K$, $3\mid K$, and
\[
 Q\equiv1\pmod{b^2},\qquad Q\equiv1\pmod3,\qquad \gcd(b,Q)=1.
\]
On every signed assignment $3\mid A$, $C\equiv-1\pmod3$, so $D\equiv0$
modulo3 and $a_1\equiv0\pmod3$. In contrast $r=bQ\equiv1\pmod3$ and
$\delta\equiv0\pmod3$, so $a_0\equiv1\pmod3$. Thus
\emph{$a_0\ne a_1$ unconditionally}, before any positivity argument.

The odd integer $S=2A-5$ is nonzero. The polynomial
$\pi(v)=(v^3-v-1)(v^2+23)$ has no rational root; hence
$\Theta=S^5\pi(\rho/S)\ne0$. A zero of
\eqref{eq:free-C-final} therefore has $\Sigma\ne0$. With $q=Z/\Sigma$,
the equation gives $(q^2-a_0-a_1)^2=4a_0a_1$.
Since $a_0\ne a_1$, $q\ne0$. The rationals
\[
 u=\frac{q^2+a_0-a_1}{2q},\qquad v=q-u
\]
have squares $a_0,a_1$. Rational-square-root integrality makes $u,v$
integers. Thus $a_0,a_1$ are squares and $\Sigma\mid\Theta$.
A prime-valuation comparison in the homogeneous quintic yields $S\mid T$:
if $v_p(T)<v_p(S)$, the two factors of $\Theta$ have valuations
$3v_p(T)$ and $2v_p(T)$, contradicting $S^5\mid\Theta$.

The factorization $a_0=r(1+\delta r z^2)$ has two coprime nonzero factors.
Since $r\equiv1\pmod3$, the signed-square possibility makes $r$ and the
bracket positive squares. Applying the same argument to $r=bQ$ gives
$b,Q$ positive squares, hence $b\ge4$. If $\Delta\le0$, then
$\delta=9\Delta-3\le-3$, making the bracket negative. Thus
$\Delta>0$. Its defining inequality gives $K\ne0$ and $0<g<2\mathcal C$.
If $l=0$ then $Q=1-K^2\le0$, a contradiction. Hence $l\ne0$;
$f_{\rm old}=2t+1$ is already nonzero.

The source bounds now give $X>3b$ and $Y>256$. Since $V+1$ is odd and
nonzero,
\[
 |A|=2|l|XY|V+1|\ge2XY>2B_1+2.
\]
All hypotheses of the strengthened lemma hold: $6\mid A$,
$B_1\equiv3\pmod4$, $C\equiv5\pmod6$, and $\xi\ne0$.
It gives $C=u_{B_1}(A)$. The recurrence modulo $A-2$ then gives
$(A-2)\mid C-B_1$ \emph{as a consequence}. Restore the proof-only integer
$h=(C-B_1)/(A-2)$ and apply the unchanged reverse binomial bridge.
The calibrated source decodes a positive source zero and hence a finite
ZFC refutation. No sign of an arbitrary supplied auxiliary variable was
silently assumed.

\subsection{Completeness}

Start from a positive source zero and its canonical arithmetic witnesses.
They have $C=u_{B_1}(A)$, $6\mid A$, and $B_1\equiv3\pmod4$.
Thus $C$ is odd and $C\equiv-1\pmod3$, giving the integer
$c_0=(C+1)/6$. The old constructive congruences still supply $k_0$.
Here $b,Q$ are positive squares and $\Delta>0$.
The positive integer
\[
 \delta=3(3\Delta-1)
\]
has $3$-adic valuation exactly1, and is therefore nonsquare; the same holds
for $\delta r$ since $r=bQ$ is a positive square.
Pell's theorem supplies positive $u,z$ with
$u^2-\delta r z^2=1$, making $a_0=ru^2$ square.

The modified Lucas interpolation theorem first supplies a positive $\xi$
divisible by this $z$, and then permits every residue class of $y$ modulo
that fixed $\xi$ while retaining $DFI$ square. Put $x=\xi/z$.
Since $S\mid T$, write $v_0=T/S$. The quintic $\pi$ has a root modulo every
positive integer, so prescribe $y$ with $\pi(v_0+y)\equiv0\pmod\xi$.
Then $\Sigma\mid\Theta$. Taking
\[
 m=\Theta/\Sigma+\sqrt{a_0}+\sqrt{a_1}
\]
gives an integer solution of \eqref{eq:free-C-final}.
Thus the construction is existence-equivalent to the fixed source, and
negating the two existence statements gives the required equivalence with
$\operatorname{Con}(\mathrm{ZFC})$.

\subsection{Exact degree and verification boundary}

Let $a=\deg A$. Since $C$ now has degree1,
\begin{align*}
 \deg D&=2a+2,& \deg E&=2a+6,& \deg F&=6a+12,\\
 \deg G&=7a+12,&\deg H&=8a+19,&\deg I&=30a+62,\\
 \deg a_1&=38a+76,&\deg\Sigma&=5a+2.
\end{align*}
For the included compatible source, the guard and center have strictly
lower degree than the dominant $\Sigma^4a_1^2$ term. Consequently
\[
 \deg F_{10}=4(5a+2)+2(38a+76)=96a+160.
\]
Exact nonzero diagonal leading coefficients are checked with the actual
fixed tuple modulo three primes. They certify that the upper degree is
attained; they are not sampled estimates of an unproved upper bound.
Separate checks authenticate coefficient bytes, exact source extraction,
full modular formulas, signed residue separation and live circuit topology.
Finite checks do not replace the unbounded mathematical proof or its
stated arithmetic/compiler dependencies.
